\begin{figure}[htbp]
    \centering
    \begin{tikzpicture}[
        % Compact dimensions to ensure a clean horizontal fit on a single page
        block/.style={
            draw, rectangle, 
            minimum height=1.1cm, 
            minimum width=1.8cm, 
            rounded corners=3pt, 
            line width=0.8pt, 
            align=center, 
            fill=white, 
            font=\scriptsize\sffamily\bfseries
        },
        nobox/.style={
            rectangle, 
            minimum height=1.1cm, 
            minimum width=1.8cm, 
            align=center, 
            font=\scriptsize\sffamily\bfseries
        },
        arrow/.style={
            -{Stealth[scale=1.2]}, 
            line width=0.8pt
        }
    ]
    
    % ==================== LINEAR SIMULATION PIPELINE ====================
    \node (source) [nobox] at (0,0) {Source\\Signal};
    \node (delay)  [block, right=0.30cm of source] {Time\\Delay};
    \node (atten)  [block, right=0.30cm of delay]  {Distance\\Attenuation};
    \node (noise)  [block, right=0.30cm of atten]  {Noise \&\\Delayed Echoes};
    \node (shadow) [block, right=0.30cm of noise]  {Binaural\\Head Shadow};
    \node (spec)   [block, right=0.30cm of shadow] {Elevation\\Filtering};
    \node (wave)   [nobox, right=0.30cm of spec]   {Received\\Waveform};

    % ==================== CONNECTIONS ====================
    \draw [arrow] (source) -- (delay);
    \draw [arrow] (delay)  -- (atten);
    \draw [arrow] (atten)  -- (noise);
    \draw [arrow] (noise)  -- (shadow);
    \draw [arrow] (shadow) -- (spec);
    \draw [arrow] (spec)   -- (wave);

    \end{tikzpicture}
    \caption{Signal transformation pipeline used to generate a simulated received waveform. Environmental noise and delayed echo components are inserted before binaural head-shadow and elevation filtering, so these disturbances can distort the cues used by later pathway models.}
    \label{fig:monaural_simulation_pipeline}
\end{figure}
